\documentclass[a4paper,12pt]{article}
\usepackage[a4paper, margin=1in]{geometry}  % You can change margin=1in to desired values
\setcounter{secnumdepth}{3}
\usepackage{comment}
\usepackage[usenames,dvipsnames,table]{xcolor}
\usepackage{pdfcolmk} 	% to correct pdflatex colour mistakes
\usepackage{cancel} %for strinking through terms which cancel oUa_1
\usepackage{fancyhdr}
\usepackage{placeins}
\usepackage{soul}
\usepackage{graphicx}
\usepackage{caption}
\usepackage{subcaption}
\usepackage{float}
%\usepackage{cite}
\usepackage{bm}
\usepackage{rotating}
\usepackage[utf8]{inputenc}
\usepackage{amsmath}
\DeclareMathOperator{\sign}{sign}
%\usepackage{natbib}%[aUa_1horyear,sort]
\usepackage{latexsym}
\usepackage{amssymb}
\usepackage{setspace}
\usepackage{amsmath}
\usepackage{framed} \definecolor{shadecolor}{cmyk}{0 0 0.2 0}
\usepackage{amsthm}
\newtheorem{definition}{Definition}
\newtheorem{proposition}{Proposition}
\newtheorem{theorem}{Theorem}
\newtheorem{example}{Example}
\newtheorem{corollary}{Corollary}
\newtheorem{lemma}{Lemma}
\usepackage{qtree}
\usepackage{makecell}
\usepackage{listings}

\lstset{
  language=Python,
  basicstyle=\ttfamily\small,
  keywordstyle=\color{blue},
  commentstyle=\color{gray},
  stringstyle=\color{orange},
  showstringspaces=false,
  breaklines=true,
  frame=single
}
\usepackage{setspace}
\doublespacing
\usepackage{endnotes}

\let\footnote=\endnote

\usepackage[colorlinks,urlcolor=blue,citecolor=blue,linkcolor=blue,pdfstartview=FitH,pdfpagemode=UseNone]{hyperref}
\usepackage{rotating}
\usepackage{authblk}
\usepackage{orcidlink}
\usepackage[normalem]{ulem}
\newcommand{\del}[1]{\sout{#1}}
\usepackage{epigraph}
%\usepackage{natbib}
%\setcitestyle{square}

\usepackage[style=verbose, citestyle=numeric, bibstyle=numeric]{biblatex}
\addbibresource{../1refs_revision.bib}

\sloppy
\usepackage{tabularx}
\newcommand{\juergen}[1]{\textcolor{cyan}{#1}}
\newcommand{\ignacio}[1]{\textcolor{red}{#1}}
\newcommand{\new}[1]{\textcolor{magenta}{#1}}

% --- reply-letter macros ---------------------------------------------------
% \reply{}  our text to the reviewer, set in blue
% \open{}   a placeholder for an item still being drafted; deliberately loud,
%           so that no unfinished reply can be shipped by accident.
% \rlabel{}  a block label (Comment. / Response. / Changes.) set on its own
%           line, so that the text beneath it always starts a fresh paragraph.
\newcommand{\reply}[1]{\textcolor{blue}{#1}}
\newcommand{\open}[1]{\textcolor{red}{\textbf{[OPEN -- TO BE COMPLETED]}\ #1}}
\newcommand{\rlabel}[1]{\par\medskip\noindent\textbf{#1}\par\nobreak}
% ---------------------------------------------------------------------------

\usepackage{lmodern}
\usepackage{graphicx}

\usepackage{tikz}
\usetikzlibrary{shapes,arrows,positioning, arrows.meta, quotes}
\usetikzlibrary{decorations.pathreplacing,shapes.misc,bending}
\usepackage{pgfplots}

\tikzset{decision/.style={diamond, draw, fill=blue!20, text width=4.5em, text badly centered, inner sep=0pt}}
\tikzset{block/.style={rectangle, draw, fill=blue!20, text width=5em, text centered, rounded corners,
 minimum width=3.5cm}}
\tikzset{line/.style={draw, -latex}}
\tikzstyle{process} = [rectangle, minimum width=3cm, minimum height=1cm, text centered, draw=black, fill=orange!30]
\tikzstyle{io} = [trapezium, trapezium left angle=70,minimum width=3cm, minimum height=1cm, trapezium right angle=110, text centered, draw=black, fill=blue!30]

\title{Vaccination Policy under Model Uncertainty: Can the Needs of the Few Outweigh the Needs of the Many? }


\author{}


\allowdisplaybreaks
\begin{document}

\begin{center}
{\Large\bf Response to the Reviewers}\\[0.6em]
{\large Vaccination Policy under Model Uncertainty:\\ Can the Needs of the Few Outweigh the Needs of the Many?}
\end{center}

\bigskip

\noindent We are sincerely grateful to both reviewers for the care and the
generosity of their reports. Reading them, we realised that several of the
points raised were not small matters of presentation but places where the
manuscript was either imprecise or simply wrong, and we have been glad of the
chance to put them right. We are particularly grateful for the references we
were pointed to, and for a set of historical counter-examples that sharpened
our thinking about the central assumption of the paper.

\medskip

\noindent Below we take each comment in turn. For every comment we give the
reviewer's words in full, then our reply, then the specific changes made to
the manuscript. Our text is set in {\color{blue}blue}. All section numbers
refer to the revised manuscript, and all changes are marked in colour there,
so that each item can be located directly.

\medskip

\noindent \open{Two items are still being drafted and are flagged in red
below: Reviewer 1's Comment 5 and Reviewer 2's Comment 1. This notice and both
placeholders must be resolved before the letter is submitted.}

\bigskip

\section*{Summary of the changes}

\noindent\begin{tabularx}{\textwidth}{|l|X|l|}
\hline
{\bf Comment} & {\bf Point raised} & {\bf Revised ms.}\\\hline
R1 $\cdot$ C1 & Engage the selective-lockdown ethics literature & \S2.4\\\hline
R1 $\cdot$ C2 & Engage the modelling-humility and uncertainty-communication literature & \S5\\\hline
R1 $\cdot$ C3 & Ground the confidence claim in the Simplifications section & \S2.5\\\hline
R1 $\cdot$ C4 & Clarify what is negative about ``negative herd immunity'' & \S3.3\\\hline
R1 $\cdot$ C5 & Specify which notion of herd immunity is in use & \open{}\\\hline
R1 $\cdot$ C6 & Introduce the block quotations properly & \S1, \S2.1\\\hline
R2 $\cdot$ C1 & The virulence-decline premise & \S2.1, \S3.2\\\hline
R2 $\cdot$ C2 & The claim about tetanus in the introduction & \S1\\\hline
\end{tabularx}

\bigskip

\section*{Reviewer 1}

\subsection*{Overall comment}

\rlabel{Comment.}
This is a very good paper which proposes an alternative model for infectious disease spread under conditions of uncertainty. The proposed model yields conclusions and supports recommendations substantially different from those normally derived from models, especially those used during the pandemic. By proposing a plausible model (that is, one based on plausible, albeit questionable, assumptions about disease spread, human behaviour, virulence changes etc) and suggesting that selective, rather than population-wide, vaccination could result in fewer deaths given certain assumptions, the paper emphasises the uncertainty and speculative nature of many epidemiological models used for policy recommendation, thus making a case for (what I take to be) much needed epistemic humility in model-based epidemiology and subsequent policy recommendations.

The paper is clearly written and well structured and I think it would make a nice contribution to the literature on pandemic/epidemic preparedness and on model-based policy making.

For these reasons, I recommend that it is published, though I have some suggestions for improvements and I should point out that my expertise is in ethics/philosophy only. I cannot therefore assess the technical aspects of the proposed model.

\rlabel{Response.}
\reply{We thank the reviewer warmly for this generous assessment and for the
recommendation to publish. We are pleased that the argumentative role of the
model came across as intended: the paper is not a prediction about any
particular pandemic, but a constructive counter-example meant to show how much
weight contestable modelling assumptions can carry. We have adopted all six
suggestions below. Four of them (C1, C2, C3 and C6) are straightforward
additions or clarifications; C4 led us to rewrite the passage the reviewer
found unclear; C5 raises a question about the notion of herd immunity that we
address separately.}

\subsection*{Comment 1: Engage the selective-lockdown ethics literature}

\rlabel{Comment.}
About the ethics of selective vaccination, the authors rightly say that additional normative premises are needed to draw ethical implications and policy recommendation from a model predicting fewer deaths with selective vaccination. The literature on selective lockdowns, which were defended ethically on the basis of similar considerations (creating immunity among the non-vulnerable by letting them free to live their normal lives while temporarily shielding the elderly and vulnerable) might offer some insights that the authors might want to engage with or at least refer to. See e.g. Savulescu J, Cameron J. Why lockdown of the elderly is not ageist and why levelling down equality is wrong. Journal of Medical Ethics 2020;46:717-721.

\rlabel{Response.}
\reply{We are grateful for this pointer. The parallel the reviewer draws is
close: shielding the vulnerable while allowing immunity to build among the
non-vulnerable is structurally very near to our Strategy 2, and the normative
premises that were used to defend selective lockdowns are of the kind our
paper says are needed here. We have added a reference to that literature at
the point where the manuscript states that such further premises are required.
We say more about the extent of our engagement in our reply to Comment 2.}

\rlabel{Changes.}
\reply{In \S2.4 (\emph{Limits of Ethical Inference from the Model}), the closing
sentence of the third paragraph now ends: ``\dots\ require additional normative
premises that lie outside its scope, such claims are, for instance, discussed
in the literature on selective lockdowns'', with references to Savulescu and
Cameron (2020), Cameron et al.\ (2021) and Giubilini et al.\ (2021).}

\subsection*{Comment 2: Engage the modelling-humility and uncertainty-communication literature}

\rlabel{Comment.}
About the uncertainty around models, their susceptibility to contested assumption, and the importance of epistemic humility that these factors, many of the recommendations put forward by the authors in the conclusion align with the recommendations in Saltelli A, et al 2020. Five ways to ensure that models serve society: a manifesto. Nature. Jun;582(7813):482-484. The authors might want to engage with or refer to those. And about the importance of communicating uncertainty around public health policy, the authors might want to see Giubilini, A et al 2025, ``Expertise, Disagreement, and Trust in Vaccine Science and Policy. The Importance of Transparency in a World of Experts''. Diametros 22 (82): 7-27.

\rlabel{Response.}
\reply{Many thanks for drawing our attention to these papers, which are indeed
close to our conclusions. Both are now cited. We should be candid about how far
we have taken this and the previous suggestion. A proper engagement with either
body of work would deserve several pages, and we judged that it would pull the
manuscript away from its own line of argument. We have therefore opted for
citation rather than discussion, and we would of course be glad to expand if
the reviewer thinks the paper would be better for it.}

\rlabel{Changes.}
\reply{In \S5 (\emph{Conclusions}), Saltelli et al.\ (2020) is now cited at the
end of the paragraph on uncertainty about what a proper epidemiological model
should look like. Giubilini et al.\ (2025) has been added to the references
supporting the recommendation of good uncertainty communication in the
penultimate paragraph.}

\subsection*{Comment 3: Ground the confidence claim in the Simplifications section}

\rlabel{Comment.}
The section on Simplifications of the model might benefit from further clarifications. In particular, the authors write at the end of that section that ``Despite all these simplifications, we are confident that the main qualitative conclusions our model licences can also be drawn from modified models and/or evaluation metrics''. Where does that confidence come from? Why would the reader trust your confidence? Surely, appeal to your own sense of confidence by itself does not provide sufficient support.

\rlabel{Response.}
\reply{This is a fair hit, and we accept it. An appeal to our own confidence
does not, on its own, provide support, and the sentence as written asked the
reader to take our word for something we had not argued. We have added the
reasons behind the claim and, equally importantly, we now say plainly how
limited the support those reasons provide is. The honest position is this: we
did not run the further simulations, so we cannot report them. What we can
offer is the heuristic that small changes to a system rarely revolutionise its
behaviour, and the observation that more complex simulations tend to be less
predictable rather than more, which if anything makes the choice between
vaccination strategies harder rather than easier. We prefer to state that and
leave readers to draw their own conclusions.}

\rlabel{Changes.}
\reply{In \S2.5 (\emph{Simplifications}), the closing paragraph has been
extended. After the sentence the reviewer quotes, it now continues: ``Since we
did not carry out further simulations, we cannot report them and have to rely
on heuristics for predicting their results. Our confidence is partly supported
by robustness heuristics that small changes to a system rarely revolutionise
system behaviour. Furthermore, more complex simulations normally lead to less
predictable outcomes. Hence, complicating our model is, in our minds, likely to
make it even harder to predict which vaccination strategy is preferable.
Ultimately, we invite readers to draw their own inferences and conclusions;
these may be supported by further modelling efforts.''}

\subsection*{Comment 4: Clarify what is negative about ``negative herd immunity''}

\rlabel{Comment.}
The notion of `negative herd immunity' the authors introduce at p.\ 14 might need to be better explained. The authors write ``This situation can be understood as a sort of negative herd immunity effect, where herd immunity may have negative consequences for those who remain susceptible''. It's not very clear why the implications are negative. I suppose herd immunity is good in any case for those susceptible as, by definition, it means a significantly reduced risk of infection (compared to a situation of non-herd immunity). It is not clear to me that `herd immunity' itself is what is negative, as opposed to the way in which it is achieved. If I understand well, on this model, vaccine-induced herd immunity would be less good than collective immunity obtained by allowing natural infection among the non-vulnerable, as the latter is more likely than the former to reduce virulence. But that simply means that the latter is preferable to the former (though neither is negative) and that what is more or less preferable is the way in which herd immunity is achieved, rather than the outcome itself.

\rlabel{Response.}
\reply{We thank the reviewer for pressing on this, and we accept the
reconstruction offered. It is not herd immunity that is bad in our model. What
matters is the route by which collective immunity is reached, and what the
model tracks is a change of sign over time rather than a bad outcome. Reducing
transmission early does protect the population, but in our model it also keeps
the virus circulating in a smaller and better-protected pool, so that it ages
more slowly. Over the full run of an epidemic, the early benefit can be
outweighed, and more agents die. The wording we had borrowed from Luyten et
al.\ (2011) obscured this, since their negative herd immunity effect names a
different mechanism, the shifting of severe illness to a later age. We have
rewritten the passage so that the temporal structure of the effect is explicit
and so that our use of the term is clearly distinguished from theirs.}

\rlabel{Changes.}
\reply{In \S3.3 (\emph{Interpretation}), the clause the reviewer quotes,
``where herd immunity may have negative consequences for those who remain
susceptible'', has been deleted. The sentence now reads: ``This situation can be
understood as a sort of negative herd immunity effect, an initial reduction of
the transmission rate grants an initial positive immunity effect to the
population (herd immunity) that, as time progresses, turns into a negative
overall effect causing more agents to die.'' The contrast with the age-related
mechanism of Luyten et al.\ is retained at the end of the same paragraph.}

\subsection*{Comment 5: Specify which notion of herd immunity is in use}

\rlabel{Comment.}
Moreover, `herd immunity' can mean different things (from a dynamic equilibrium that allows for periodical waves of the virus, as per Sunetra Gupta's understanding, to simply very low risk that a susceptible and an infectious person come into contact). It would be good to clarify which understanding of herd immunity is used here, as that seems relevant to an epidemilogical assessment of the model.

\rlabel{Response.}
\open{This reply is still being drafted. The substantive question is which of
the two notions the reviewer distinguishes our SISD model actually realises.
Since the model has no permanent removed compartment, recovered agents return
to the susceptible pool carrying accumulated immunity, so the operative notion
appears to be closer to a dynamic equilibrium than to a threshold that is
crossed once. This has to be settled and stated in the manuscript before the
reply can claim it.}

\rlabel{Changes.}
\open{None yet.}

\subsection*{Comment 6: Introduce the block quotations properly}

\rlabel{Comment.}
Small comment: There are citations that are added to the text without context or being introduced properly. For instance it is unclear where the citation at pp.5-6 comes from and how it relates to the discussion there, and same for citation at p.\ 8. Citations should be properly introduced in a way that makes them fit with the flow of the text.

\rlabel{Response.}
\reply{We fully agree, and we are grateful for the observation. The passages in
question are block quotations that we use to support our own reasoning, but we
had dropped them into the text without saying whose words they are or why they
appear at that point. Each is now introduced by name.}

\rlabel{Changes.}
\reply{In \S1 (\emph{Introduction}), the quotation from Br\"ussow is now
preceded by ``Br\"ussow wrote during the pandemic:''. In \S2.1 (\emph{Main
Components of the Model}), under the heading on viral age, the quotation from
Kun et al.\ is now preceded by ``In the words of Kun et al.:''.}

\bigskip

\section*{Reviewer 2}

\subsection*{Overall comment}

\rlabel{Comment.}
Thank you for asking me to review this paper. The more general points in the paper about the limits of modelling and of assumptions that everyone must be vaccinated are important points.

\rlabel{Response.}
\reply{We thank the reviewer for the time given to the manuscript, and for
granting the general points about the limits of modelling. The central
objection that follows is, in our view, the most serious challenge the paper
has received, and we take it up in detail below.}

\subsection*{Comment 1: The virulence-decline premise}

\rlabel{Comment.}
However, the main thrust of the paper depends on a doubtful empirical claim, that pandemic viruses decline in virulence. While many people might believe this claim, and while there might be evidence for effects of this kind, such effects are clearly dwarfed by the effects (acting in the same direction on average disease severity) of acquired immunity, whether via vaccination or (as was more common in historic pandemics) previous infection.

Consider these counter-examples: (1) Hong Kong \& omicron - a very high death rate was observed due to a lack of immunity due to low prior infection rates and low vaccination among older adults; in other words, low population immunity was a much bigger factor than any decline in the intrinsic virulence of omicron. (2) remote communities of the pacific or the arctic where epidemics of 1918 flu were typically late but nonetheless severe. More generally (3) the high mortality from smallpox and measles when introduced to immuno-naive populations via colonisation, suggesting again that a lack of prior immunity was a much bigger factor than decline in virulence from centuries of endemic transmission in Europe.

These difficulties with the core assumption make it difficult to evaluate any ethical implications that follow.

\rlabel{Response.}
\open{This reply is still being drafted. The material available for it
includes: that our model represents acquired immunity as a separate channel
from viral ageing, so the two effects are not conflated; that the
argumentative role of the paper is modal rather than predictive, so the premise
needs to be possible and not unreasonable rather than typical; that we found
parameter settings with a small viral ageing effect in which Strategy 2 still
outperforms Strategy 1; and an honest concession that the size of the
virulence-decline effect in real pandemics is contested. Whether the reply also
reports a new robustness check remains to be decided.}

\rlabel{Changes.}
\reply{Two changes prompted by this comment have already been made to the
manuscript. In \S2.1 (\emph{Main Components of the Model}), a new paragraph
follows the description of viral age. It cites all three of the reviewer's
counter-examples, namely the Omicron BA.2 variant in Hong Kong, the isolated
Pacific islands after the 1918--1919 influenza pandemic, and immuno-naive
populations encountering pathogens from European colonisers, and it grants that
``virulence decline can be small and dwarfed by acquired immunity or even
non-existent''. It then sets against this the evidence that immunity to
Covid-19 waned over time and fell sharply after Omicron. In \S3.2, a sentence
now directs the reader to those parameter settings with a small viral age
effect that correspond to the cases the reviewer raises.}

\subsection*{Comment 2: The claim about tetanus in the introduction}

\rlabel{Comment.}
The introduction also contains the incorrect claim that tetanus has been virtually eradicated via herd immunity - this is not true, at least not insofar as ``herd immunity'' means something like reduced population transmission arising from the presence of transmission-reducing immunity in large numbers of individuals, because tetanus is not transmitted between people.

\rlabel{Response.}
\reply{The reviewer is right, and we are grateful for the correction. Tetanus
is contracted from environmental spores and is not transmitted between people,
so it cannot be controlled by herd immunity and should not have been offered as
an example of it. The opening sentence has been corrected.}

\rlabel{Changes.}
\reply{In \S1 (\emph{Introduction}), the opening sentence previously read
``\dots\ such as the eradication of smallpox and the virtual eradication of
tetanus and polio by achieving herd immunity''. It now reads ``\dots\ such as
the eradication of smallpox and the virtual eradication of tetanus (among
vaccinated patients) and polio''. The appeal to herd immunity has been dropped
from the sentence.}

\bigskip

\section*{Further changes}

\noindent\reply{Beyond the points raised by the reviewers, we have made two
editorial passes over the manuscript. Em dashes have been replaced by
parentheses throughout, and multiple citations within a single bracket have
been ordered alphabetically. In \S3.5 (\emph{Worst Cases}) we have also added a
short remark that decision makers are likely to be reluctant to bear the costs
should selective vaccination backfire, which makes worst-case outcomes of
practical as well as theoretical interest.}

\bigskip

\noindent\reply{We thank both reviewers again for their attention to the
manuscript, which we believe is considerably stronger for it.}

\end{document}
